\section{System model, adversary model, and notation}
\label{sec:model}\label{sec:02-system-model}

Every property this paper states later --- about block production (\S\ref{sec:consensus}), about
revenue-funded compensation (\S\ref{sec:pnr}), about delegated authority (\S\ref{sec:agent-identity}),
about private payment (\S\ref{sec:private-renewal}), about bridging (\S\ref{sec:bridge}) and about
governance (\S\ref{sec:governance}) --- is a statement about the model fixed here. This section names
the participants and what each can do; states the two different consistency models the system runs
(one for the chain, one for the orchestration layer) and says why the difference is the most
important structural fact about \Flux; enumerates the trust boundaries as they exist in the deployed
code, including the ones a reader would not expect; gives the adversary an explicit capital budget
denominated in \flux and prices what each level of capital buys against the measured network; and
records eleven numbered assumptions, \textbf{A1}--\textbf{A11}, that every later section cites by
number. The assumptions are stated neutrally. Several of them foreclose properties a reader might
otherwise assume hold; that is the point of writing them down.

Unless a claim is tagged otherwise, this section is \shipped and each claim carries either a source
citation of the form \srccode{repo/\allowbreak{}file}{line} or a measurement citation of the form
\srcmeasured{endpoint}{date}.

\subsection{Participants}
\label{sec:model:participants}

\begin{definition}[Node operator]
\label{def:operator}
\shipped A \emph{node operator} is a participant that controls, for some tier
$t \in \tiers = \{\tC,\tN,\tS\}$: (i) an unspent transparent output of exactly $\coll{t}$, with
$\coll{\tC} = 1{,}000\flux$, $\coll{\tN} = 12{,}500\flux$, $\coll{\tS} = 40{,}000\flux$
\srccode{fluxd/\allowbreak{}src/\allowbreak{}fluxnode/\allowbreak{}fluxnode.h}{57--63}; (ii) a node signing key, used to sign the
registration and confirmation transactions and to sign \PoN{} (Proof of Node) blocks; and (iii) a payment address,
which need not equal either of the above.

\emph{Capabilities.} Broadcast a \texttt{FLUXNODE\_START\_TX} spending a collateral output matured
by at least $100$ blocks \srccode{fluxd/\allowbreak{}src/\allowbreak{}coins.cpp}{630--686}; broadcast a benchmark-signed
confirmation transaction to enter the paid set \srccode{fluxd/\allowbreak{}src/\allowbreak{}fluxnode/\allowbreak{}activefluxnode.cpp}{18--74};
re-confirm at least every $640$ blocks or be removed \srccode{fluxd/\allowbreak{}src/\allowbreak{}fluxnode/\allowbreak{}fluxnode.h}{31--34};
run one \FluxOS{} instance that participates in the orchestration overlay and self-selects
applications; and, in any slot $\tslot$ for which its collateral outpoint is eligible, produce and
sign one block.

\emph{Economic position.} Every block pays exactly one node of each tier a reward of
$\subsidypon(\hgt)\cdot\tbase{t}/14$, with tier bases $\tbase{\tC} = 1.0$, $\tbase{\tN} = 3.5$,
$\tbase{\tS} = 9.0$ against $\subsidypon(\hpon) = 14\flux$
\srccode{fluxd/\allowbreak{}src/\allowbreak{}main.cpp}{2886--2939}. Because the per-tier payment queue is ordered by
longest-time-since-last-paid \srccode{fluxd/\allowbreak{}src/\allowbreak{}fluxnode/\allowbreak{}fluxnode.h}{313--328}, the rotation length
$\rot{t}$ equals the tier's node count $\ncount{t}$. Collateral is locked but is not destroyed by any
lifecycle transition: the three removal conditions are collateral spent, failure to re-confirm, and
collateral-amount mismatch during the migration windows
\srccode{fluxd/\allowbreak{}src/\allowbreak{}fluxnode/\allowbreak{}fluxnode.cpp}{231--290}. There is therefore no slashing, and the cost of
misbehaviour to an operator is the opportunity cost of locked capital, not a forfeited bond.

\emph{Measured population.} $\ncount{\mathrm{tot}} = 6{,}489$ confirmed nodes, of which
$\ncount{\tC} = 3{,}247$, $\ncount{\tN} = 1{,}615$, $\ncount{\tS} = 1{,}627$; total collateral
$88{,}514{,}500\flux$ \srcmeasured{api.runonflux.io/\allowbreak{}daemon/\allowbreak{}listzelnodes}{2026-09-03}.
\end{definition}

\begin{definition}[Tenant]
\label{def:tenant}
\shipped A \emph{tenant} is a participant that controls a FluxID (a base58 P2PKH-style address) and a
spending key over transparent \flux. It holds no collateral, has no consensus role, and is
identified to the orchestration layer only by an ECDSA signature over a login phrase or over an
application specification \srccode{flux/\allowbreak{}ZelBack/\allowbreak{}src/\allowbreak{}services/\allowbreak{}verificationHelper.js}{17--48}.

\emph{Capabilities.} Register, update and remove a globally-registered application; pay the
computed price to a fixed transparent sink address; start, stop and exec into its own containers via
the \texttt{appownerabove} privilege \srccode{flux/\allowbreak{}ZelBack/\allowbreak{}src/\allowbreak{}routes.js}{1196--1265}.

\emph{Sub-types.} A \emph{human} tenant reviews and signs on a device it controls. An
\emph{autonomous} tenant --- an agent $\agent$ acting for a principal $\principal$ --- differs in
exactly one respect that matters to this paper: it must be able to act without a human in the loop,
and the deployed stack already permits that through a bearer token whose authority is bounded by
action and unbounded by value (\S\ref{sec:agent-identity}). No delegated-spending primitive with a
budget $\budget$ or a scope $\scopeset$ exists in the deployed system.
\end{definition}

\begin{definition}[The Foundation as a distinguished participant]
\label{def:foundation}
\shipped Write $\mathcal{F}$ for the participant defined \emph{extensionally}, as the holder of the
following key sets and the operator of the following infrastructure. $\mathcal{F}$ is distinguished
because several of its capabilities cannot be acquired with any amount of \flux.

\begin{itemize}
  \item[\textbf{K1.}] The $4$ hardcoded emergency public keys; any $2$ of them suffice to author a block that
        bypasses \PoN{} eligibility and per-node signature verification
        \srccode{fluxd/\allowbreak{}src/\allowbreak{}chainparams.cpp}{242--253}, \srccode{fluxd/\allowbreak{}src/\allowbreak{}emergencyblock.cpp}{56--150}.
  \item[\textbf{K2.}] The $5$-entry benchmark-signing allowlist \texttt{vec\allowbreak{}Benchmarking\allowbreak{}Public\allowbreak{}Keys}, each entry
        carrying a UNIX-timestamp validity start, against which every node's benchmark signature is
        verified by consensus \srccode{fluxd/\allowbreak{}src/\allowbreak{}chainparams.cpp}{233--240}.
  \item[\textbf{K3.}] The dev-fund recipient address, to which a consensus-enforced minimum coinbase output of
        $\subsidypon(\hgt) - \sum_{t\in\tiers}\subsidypon(\hgt)\tbase{t}/14$ --- i.e.\ $0.5/14 =
        3.57\%$ of \PoN{} emission at activation --- must be paid
        \srccode{fluxd/\allowbreak{}src/\allowbreak{}main.cpp}{2877--2884}, enforced at \srccode{fluxd/\allowbreak{}src/\allowbreak{}main.cpp}{1231--1249}.
  \item[\textbf{K4.}] Two FluxIDs holding the \texttt{fluxteam} privilege on \emph{every} node's \FluxOS{} API,
        conferring \texttt{adminandfluxteam} rights including broadcast injection, daemon
        start/restart, and application updates on behalf of an owner
        \srccode{flux/\allowbreak{}ZelBack/\allowbreak{}src/\allowbreak{}services/\allowbreak{}verificationHelperUtils.js}{96--156},
        \srccode{flux/\allowbreak{}ZelBack/\allowbreak{}src/\allowbreak{}routes.js}{1386--1441}.
  \item[\textbf{K5.}] Merge access to the policy repositories whose documents every node fetches and enforces
        (\texttt{fluxos-\allowbreak{}network-\allowbreak{}policy}, and the \texttt{helpers/} mirror in \texttt{RunOnFlux/flux})
        \srccode{fluxos-network-policy/\allowbreak{}.github/\allowbreak{}CODEOWNERS}{1--15},
        \srccode{flux/\allowbreak{}ZelBack/\allowbreak{}src/\allowbreak{}services/\allowbreak{}appSecurity/\allowbreak{}imageManager.js}{179--195}.
  \item[\textbf{K6.}] The UEFI Secure Boot Platform Key whose SHA-256 is pinned in the \ArcaneOS{} enrolment
        binary; the node checks the firmware's enrolled PK against that pin
        \srccode{fluxsas/\allowbreak{}src/\allowbreak{}fes/\allowbreak{}fes.cpp}{70--81}.
  \item[\textbf{K7.}] The \fluxsas{} attestation signing material and the internal mTLS CA
        \srcprivate{\SAS{} API server}, \srccode{fluxsas/\allowbreak{}tools/\allowbreak{}gen-certs.sh}{38}.
  \item[\textbf{K8.}] The multisig that posts on-chain chain-parameter messages, which change the deployment price
        table with effect from their occurrence height
        \srccode{flux/\allowbreak{}ZelBack/\allowbreak{}src/\allowbreak{}services/\allowbreak{}utils/\allowbreak{}chainUtilities.js}{9--52}.
  \item[\textbf{K9.}] The registrar, authoritative DNS and ACME accounts for the \texttt{runonflux.io}
        name space, the $16$-host \FDM{} fleet, the $5$ PowerDNS servers, the single-host DNS gateway
        and the Arcane VIP, all inside one privately-operated \texttt{10.100.0.0/24}
        \srccode{flux-fdm-infrastructure/\allowbreak{}ansible/\allowbreak{}inventory/\allowbreak{}hosts.yaml}{1--206}.
\end{itemize}

\decided{custody is \emph{personal}, not corporate. No legal entity holds these key sets: four
named individuals --- Tadeas Kmenta, Valter Silva, Daniel Keller and Jeremy Anderson --- each hold
one, personally governed \srcintent{design intake}, matching the 3-of-4 emergency
and bridge configuration of \Cref{sec:governance}. The distinction bears on every concentration
argument in this paper: personal custody by four named people is a different accountability model
from custody by a company, and both differ from the anonymous-operator model a reader of a
decentralized-network paper might otherwise assume. No repository or endpoint read for this paper
establishes it; it rests on the author's statement.}
\end{definition}

\begin{definition}[Holder of a parallel asset]
\label{def:parallel-holder}
\shipped A \emph{parallel-asset holder} controls a token on one of ten non-\Flux{} chains that
represents \flux{}. Such a holder has no consensus role, no collateral, and no orchestration
capability. Its economic position is a claim against the operator of the incumbent \Fusion{} bridge,
not against \Flux{} consensus; \S\ref{sec:bridge} treats the resulting custody structure. The
inventory of ten chains is given in \S\ref{sec:parallel-assets}.
\end{definition}

\subsection{Network model}
\label{sec:model:network}

\begin{definition}[Chain network]
\label{def:chainnet}
\shipped The chain network is the set of \fluxd{} instances exchanging blocks and transactions over a
Bitcoin-style gossip protocol \srccode{fluxd/\allowbreak{}src/\allowbreak{}chainparams.cpp}{175--180}. We assume
\emph{partial synchrony}: there is an unknown global stabilisation time after which any message
between two correct nodes is delivered within an unknown finite bound. Before it, messages may be
delayed arbitrarily.

\PoN{} additionally requires bounded \emph{clock} skew, not merely bounded delay: the slot index is
$\tslot(t) = \lfloor (t - t_{\mathrm{genesis}})/\spacingpon \rfloor$ with
$\spacingpon = \qty{30}{\second}$ \srccode{fluxd/\allowbreak{}src/\allowbreak{}pon/\allowbreak{}pon.cpp}{88--93},
\srccode{fluxd/\allowbreak{}src/\allowbreak{}chainparams.cpp}{105}, and the only enforcement is a
$\qty{300}{\second}$ future-time tolerance on \PoN{} headers plus a strict
$\text{blocktime} > \text{prev blocktime}$ ordering
\srccode{fluxd/\allowbreak{}src/\allowbreak{}main.cpp}{5330--5378,5564--5572}.
\end{definition}

\begin{definition}[Orchestration overlay]
\label{def:overlay}
\shipped The orchestration overlay is the set of \FluxOS{} instances exchanging JSON-over-WebSocket
messages, each wrapped in an envelope
$\{\textsf{version}, \textsf{timestamp}, \textsf{pubKey}, \textsf{signature}, \textsf{data}\}$ signed
by the node key \srccode{flux/\allowbreak{}ZelBack/\allowbreak{}src/\allowbreak{}services/\allowbreak{}utils/\allowbreak{}fluxBroadcastHelper.js}{11--26}. There is no
consensus: application state lives in per-node MongoDB collections --- the \texttt{globalzelapps}
database of messages, specifications and locations that every node replicates by gossip
\srccode{flux/\allowbreak{}ZelBack/\allowbreak{}config/\allowbreak{}default.js}{68--76}.

Delivery is best-effort and threshold-gated. Duplicate suppression is a per-node message cache after
a random \qty{1}{\milli\second}--\qty{75}{\milli\second} jitter; a peer sending five malformed or
badly-signed messages in ten minutes is disconnected
\srccode{flux/\allowbreak{}ZelBack/\allowbreak{}src/\allowbreak{}services/\allowbreak{}fluxCommunication.js}{549--618}. Above chain height $1{,}751{,}250$
only the $40$-hex message hash is gossiped and the body is pulled on demand
\srccode{flux/\allowbreak{}ZelBack/\allowbreak{}src/\allowbreak{}services/\allowbreak{}fluxCommunication.js}{494--498}. Convergence state is soft: a
location record expires after \qty{7500}{\second}, a temporary message after \qty{3600}{\second}, an
installing record after \qty{900}{\second}, and a graceful-exit notification extends an affected
location by \qty{420}{\second} \srccode{flux/\allowbreak{}ZelBack/\allowbreak{}config/\allowbreak{}default.js}{311--314},
\srccode{flux/\allowbreak{}ZelBack/\allowbreak{}src/\allowbreak{}services/\allowbreak{}utils/\allowbreak{}appConstants.js}{86}.
\end{definition}

\begin{property}[Two consistency models --- the structural fact]
\label{prop:two-models}
\shipped Chain state is totally ordered and agreed by consensus. Application state --- which
applications exist, which nodes run them --- is a set of per-node databases converged by signed but
unordered gossip, with no total order and no termination guarantee. Consequently no
predicate of the form ``application $\app$ is running on node $i$ at height $\hgt$'' is decidable
from the chain.
\end{property}

\begin{proof}[Proof sketch]
Placement is opportunistic self-selection: each node independently polls for applications missing
instances and picks \emph{uniformly at random} among the candidates, subject to soft deferrals and a
broadcast collision-avoidance wait \srccode{flux/\allowbreak{}ZelBack/\allowbreak{}src/\allowbreak{}services/\allowbreak{}appLifecycle/\allowbreak{}appSpawner.js}{254--336},
\srccode{flux/\allowbreak{}ZelBack/\allowbreak{}config/\allowbreak{}default.js}{321}. The resulting placement is published as a TTL'd
location broadcast on the overlay, and nothing in the deployed code writes placement to the chain
\srccode{flux/\allowbreak{}ZelBack/\allowbreak{}config/\allowbreak{}default.js}{56--60}. The chain carries the \emph{payment} and the
registry of collateralised nodes; it does not carry the assignment of workloads to nodes. Hence any
chain-verifiable statement about placement would have to be reconstructed from off-chain gossip,
which by Definition~\ref{def:overlay} offers no agreement property.
\end{proof}

\begin{remark}[Why this is the load-bearing fact]
\label{rem:why-structural}
Four later sections turn on Property~\ref{prop:two-models}. \S\ref{sec:pnr} must place the fee pool
$\pool$ in \emph{consensus} state rather than in overlay state, because a value that determines
payment cannot be converged best-effort. \S\ref{sec:private-renewal} inherits the consequence that
the payment is the only chain-visible artefact of a deployment, which is what makes payer/workload
unlinkability hard. \S\ref{sec:bridge} cannot export any placement or capacity fact to another chain.
\S\ref{sec:governance} must site the moderation quorum on the chain, because a vote counted on the
overlay has no agreed tally. A reader who carries a single mental model of ``the \Flux{} state''
will draw wrong conclusions in all four.
\end{remark}

\paragraph{What breaks when delivery assumptions fail.}
\shipped (i) \emph{Chain partition}: competing tips are resolved by block count and first-seen rather
than by accumulated work (\textbf{A3}); reorganisations are capped at $40$ blocks
\srccode{fluxd/\allowbreak{}src/\allowbreak{}main.h}{74}, a bound that was raised to $5{,}000$ for the closed height window
$[2{,}020{,}000,\,2{,}025{,}000]$ during \PoN{} stabilisation
\srccode{fluxd/\allowbreak{}src/\allowbreak{}main.cpp}{8983--8988}. (ii) \emph{Overlay degradation}: below $12$ live peers the
application spawner is paused and below $4$ it is treated as degraded; a registration is refused
outright below $8$ outbound or $4$ inbound peers
\srccode{flux/\allowbreak{}ZelBack/\allowbreak{}config/\allowbreak{}default.js}{262--269}. (iii) \emph{Discovery outage}: with
\texttt{api.\allowbreak{}runonflux.\allowbreak{}io} unreachable, \FDM{} keeps serving its last configuration and propagates no
new information --- fail-static, with no independent source
\srccode{flux-domain-manager/\allowbreak{}src/\allowbreak{}services/\allowbreak{}flux/\allowbreak{}dataFetcher.js}{47--91}. (iv) \emph{Policy outage}: a
node that cannot fetch the policy documents keeps enforcing its last-read copy indefinitely and
across restarts, and a node that never had one declines to answer --- this is the publisher-side
contract \srcdoc{fluxos-network-policy/README.md:33--44}; the node-side fetch-and-cache path is
verified in \FluxOS{} \srccode{flux/\allowbreak{}ZelBack/\allowbreak{}src/\allowbreak{}services/\allowbreak{}appSecurity/\allowbreak{}imageManager.js}{179--195}.
(v) \emph{Hash-list outage}: if both dm-verity
allow-list endpoints are unreachable, the \ArcaneOS{} watchdog treats the node as secure and takes no
action --- fail-open \srcprivate{\SAS{} watchdog}.

\subsection{Trust boundaries}
\label{sec:model:trust}

Each row of Table~\ref{tab:trust} names something that is trusted, the participants who must trust
it, and what a failure of that trust yields. The list is not a critique; it is the set of
assumptions any later property must either use or discharge.

%% longtable, not table[H]: eleven trust boundaries do not fit one page and the
%% fixed float ran off the bottom.
{\small
\begin{longtable}{@{}p{0.045\linewidth}p{0.235\linewidth}p{0.185\linewidth}p{0.44\linewidth}@{}}
\caption{Trust boundaries in the deployed system. All rows are \shipped.}
\label{tab:trust}\\
\toprule
 & \textbf{What is trusted} & \textbf{Trusted by} & \textbf{A failure yields} \\
\midrule
\endfirsthead
\multicolumn{4}{@{}l}{\emph{(Table \ref{tab:trust} continued)}}\\
\toprule
 & \textbf{What is trusted} & \textbf{Trusted by} & \textbf{A failure yields} \\
\midrule
\endhead
\bottomrule
\endfoot
TB1 & Consensus rules in \fluxd{} \srccode{fluxd/\allowbreak{}src/\allowbreak{}main.cpp}{5330--5572} & all participants &
Divergent ledgers; reorganisation bounded at $40$ blocks. \\
TB2 & Collateral registry (\texttt{determ\_zelnodes} cache, rebuilt at start)
\srccode{fluxd/\allowbreak{}src/\allowbreak{}fluxnode/\allowbreak{}fluxnodecachedb.cpp}{49--79} & consensus and placement (and, under
\S\ref{sec:governance}, vote weight) &
Crash inconsistency between the UTXO set and the node cache; a dedicated sync marker
\srccode{fluxd/\allowbreak{}src/\allowbreak{}fluxnode/\allowbreak{}fluxnodecachedb.cpp}{17,166--171} and an
operator-invoked \texttt{rebuildfluxnodedb} repair path \srccode{fluxd/\allowbreak{}src/\allowbreak{}rpc/\allowbreak{}fluxnode.cpp}{78} exist. \\
TB3 & The \emph{local} \fluxbench{} process, invoked by \texttt{popen()}
\srccode{fluxd/\allowbreak{}src/\allowbreak{}fluxnode/\allowbreak{}benchmarks.cpp}{227--291} & the co-located \fluxd{} &
Benchmark-gating bypass: an operator's own host can produce the signature attesting that its hardware
meets the tier thresholds. Not tier forgery --- the registration must still match an on-chain
collateral amount. \\
TB4 & Benchmark maintainer allowlist \srccode{fluxd/\allowbreak{}src/\allowbreak{}chainparams.cpp}{233--240} & consensus &
Any holder of a currently-valid allowlist key can authorise benchmark signatures fleet-wide. \\
TB5 & \fluxsas{} attestation service \srcprivate{\SAS{} API server} & \fluxbench{} and
\texttt{fdm-nginx} only; \emph{no tenant path} & An attestation demonstrates possession of a value
present in every copy of the software, not that a machine booted a given image. \\
TB6 & dm-verity hash allow-list service (two Flux-operated hosts)
\srcprivate{\SAS{} watchdog} & every \ArcaneOS{} node & Whoever controls those names
defines ``known-good''; if both are unreachable the watchdog does nothing (fail-open). \\
TB7 & UEFI Platform Key pinned by SHA-256 \srccode{fluxsas/\allowbreak{}src/\allowbreak{}fes/\allowbreak{}fes.cpp}{70--81}; TPM seal against
PCR~7 only \srccode{fluxsas/\allowbreak{}src/\allowbreak{}fes/\allowbreak{}fes.cpp}{173--186} & the node's own boot path & Whoever controls
that PK defines what may boot. PCR~7 measures Secure Boot policy, so the seal opens for any
PK-signed kernel, not a specific measured image. \\
TB8 & OS update channel: release metadata and checksum files on the Flux CDN
\srccode{arcaneos-releases/\allowbreak{}publish\_release.sh}{37,65--73,98--101} & operators updating \ArcaneOS{} &
Whoever controls the CDN and metadata defines the candidate image set. \\
TB9 & \texttt{fluxos-\allowbreak{}network-\allowbreak{}policy} documents, unsigned, fetched over plain HTTPS and enforced
hourly by every node \srccode{flux/\allowbreak{}ZelBack/\allowbreak{}src/\allowbreak{}services/\allowbreak{}appSecurity/\allowbreak{}imageManager.js}{179--195},
\srccode{flux/\allowbreak{}ZelBack/\allowbreak{}src/\allowbreak{}services/\allowbreak{}serviceManager.js}{530--531}; publication contract
\srcdoc{fluxos-network-policy/README.md:30--44,320--326} & every \FluxOS{} node & Any org owner
merging to \texttt{main} can block a container image, an owner address or an application hash
network-wide, or revoke an enterprise-node grant, with effect on the next scheduled fetch:
\qty{6}{\hour} for blocked repositories and enterprise nodes, \qty{12}{\hour} for the tampering list.
No signature, no second reviewer, no staged rollout. \\
TB10 & \texttt{api.\allowbreak{}runonflux.\allowbreak{}io} as sole application-discovery source
\srccode{flux-domain-manager/\allowbreak{}src/\allowbreak{}services/\allowbreak{}flux/\allowbreak{}dataFetcher.js}{47--91} & \FDM{},
\texttt{flux-\allowbreak{}apps-\allowbreak{}dns-\allowbreak{}manager} & No independent answer to ``which applications exist and where''; the
routing plane fails static on the last-known answer. \\
TB11 & \FDM{}/DNS/certificate path: $16$ \FDM{} hosts, $5$ PowerDNS servers, one DNS gateway host,
VRRP VIP, one private \texttt{/24}
\srccode{flux-fdm-infrastructure/\allowbreak{}ansible/\allowbreak{}inventory/\allowbreak{}hosts.yaml}{1--206} & every tenant reachable by
name, every TLS client & Name resolution, HAProxy election and certificate issuance for the
application name space are controlled by one operator; a wildcard blacklist in the \FDM{} config
removes an application from that path unilaterally
\srccode{flux-domain-manager/\allowbreak{}config/\allowbreak{}appsConfig.js}{2--6}. \\
TB12 & Application payment sink addresses \srccode{flux/\allowbreak{}ZelBack/\allowbreak{}config/\allowbreak{}default.js}{241--244} &
every paying tenant & Payment is a transfer to a fixed transparent address; entitlement follows from
each node independently re-deriving the price. Underpayment is logged and the message is never
promoted, which is indistinguishable to the payer from non-propagation
\srccode{flux/\allowbreak{}ZelBack/\allowbreak{}src/\allowbreak{}services/\allowbreak{}appMessaging/\allowbreak{}messageVerifier.js}{733--746,777}. \\
TB13 & \FluxOS{} \texttt{fluxteam} privilege
\srccode{flux/\allowbreak{}ZelBack/\allowbreak{}src/\allowbreak{}services/\allowbreak{}verificationHelperUtils.js}{96--156} & every node & Two FluxIDs
can broadcast on the overlay, restart the daemon and update applications on behalf of their owners,
on every node. \\
TB14 & Chain-parameter multisig posting price messages
\srccode{flux/\allowbreak{}ZelBack/\allowbreak{}src/\allowbreak{}services/\allowbreak{}utils/\allowbreak{}chainUtilities.js}{9--52} & every tenant and node &
The deployment price schedule changes from the occurrence height of a posted message. \\
TB15 & The $2$-of-$4$ emergency key set \srccode{fluxd/\allowbreak{}src/\allowbreak{}emergencyblock.cpp}{183--191} &
all participants & Two key holders can insert a valid block at any height after \PoN{} activation,
bypassing sortition and per-node signature checks, with no cooldown, no stall detection and no
on-chain justification. \\
\end{longtable}
}

\decided{the emergency-block path is observed \emph{on-chain}: an emergency block is a block, so
its production is visible to anyone following the chain and detection does not depend on an off-chain
alerting system \srcintent{design intake}. That is the whole of the monitoring
story --- this paper does not claim a rate limiter, an alarm or a watch rota outside the chain,
because none is established and none is asserted. Observability is not the same as response, and
\Cref{sec:governance}'s move to 3-of-4 with fresh keys is the control that bounds the path.}
\notestablished{whether \ArcaneOS{} release artefacts are signature-verified beyond the published
SHA-256 sums is not established in the repositories read (TB8).}
\notestablished{the fetch interval for \texttt{vettedrepositories.\allowbreak{}json} is not stated in the policy
repository, so the \qty{6}{\hour}--\qty{12}{\hour} bound in TB9 covers four of the five documents.}
\decided{the application payment sink is held by the same four individuals as the emergency and
bridge keys above --- Tadeas Kmenta, Valter Silva, Jeremy Anderson and Daniel Keller
\srcintent{design intake}. The address itself is expected to change on migration to
SSP Enterprise; the custody does not. This is a policy fact rather than a code fact: no artefact read
for this paper establishes it, and it rests on the author's statement. It is stated because
\Cref{sec:pnr}'s settlement sends every application payment to that one address, so who holds its key
is load-bearing for the whole revenue path and not a detail a reader should have to infer.}

\subsection{The adversary}
\label{sec:model:adversary}

\begin{definition}[Adversary $\adv$]
\label{def:adversary}
$\adv$ is parameterised by the triple $(\advbudget,\ \mathcal{K},\ \mathcal{N})$:
\begin{itemize}
  \item $\advbudget \in \mathbb{R}_{\ge 0}$, a \emph{capital budget} denominated in \flux{} and
        spendable on collateral or on deployment payments. By \textbf{A11} it is never converted to
        or from fiat.
  \item $\mathcal{K} \subseteq \{\mathrm{TB1},\dots,\mathrm{TB15}\}$, the set of trust boundaries
        $\adv$ controls. Membership in $\mathcal{K}$ is \emph{not} priced in \flux{}: no budget buys
        TB9 or TB15.
  \item $\mathcal{N} \in \{\textsf{observe},\ \textsf{delay},\ \textsf{partition}\}$, network power
        over the two networks of \S\ref{sec:model:network}.
\end{itemize}
$\adv$ is computationally bounded: it cannot forge ECDSA or Ed25519 signatures under keys it does not
hold, and cannot invert SHA-256. It may deviate arbitrarily from any behaviour that is not
consensus-enforced --- in particular from the advisory rank delay of Property~\ref{prop:share}.
\end{definition}

\paragraph{What $\advbudget = 0$ buys.} \shipped Several properties in this paper fail to an
adversary with no capital at all.
\begin{itemize}
  \item[\textbf{Z1.}] \emph{Produce attestations that any consumer of \fluxsas{} accepts.} The signing seed for the
        default, \texttt{cdn}, \texttt{mesh} and \texttt{app} attestation purposes is derived from
        values compiled identically into every copy of the software; no input depends on the TPM, on
        Secure Boot state, or on any per-machine value
        \srcprivate{\SAS{} API server}. No node, no hardware and no
        collateral are required. This is why \textbf{A5} is stated as it is.
  \item[\textbf{Z2.}] \emph{Link payments to deployments.} All application payments land at fixed transparent
        addresses \srcmeasured{api.runonflux.io/\allowbreak{}apps/\allowbreak{}deploymentinformation}{2026-09-03}, and a
        payment names the application it funds; \S\ref{sec:private-renewal} measures the resulting
        anonymity set as a singleton within its block.
  \item[\textbf{Z3.}] \emph{Inject unsigned frames into the overlay.} The raw hash-gossip frames
        \texttt{message\allowbreak{}Hash\allowbreak{}Present} and \texttt{request\allowbreak{}Message\allowbreak{}Hash} bypass the signed envelope
        entirely \srccode{flux/\allowbreak{}ZelBack/\allowbreak{}src/\allowbreak{}services/\allowbreak{}fluxCommunication.js}{511--536}. The capability is
        bounded --- these frames carry a $40$-hex hash and no application data --- but it is free.
  \item[\textbf{Z4.}] \emph{Degrade.} Reducing a node's live peer count below the spawner thresholds, or causing
        confirmation loss, removes that node's applications; the enforcement is redundant across the
        fleet and therefore also amplifies a stale view
        \srccode{flux/\allowbreak{}ZelBack/\allowbreak{}src/\allowbreak{}services/\allowbreak{}appMonitoring/\allowbreak{}nodeStatusMonitor.js}{62--149}.
  \item[\textbf{Z5.}] \emph{Act as a tenant.} The price floor for a deployment is $0.01\flux$ per default
        term \srcmeasured{api.runonflux.io/\allowbreak{}apps/\allowbreak{}deploymentinformation}{2026-09-03}, so tenant-side
        capabilities are effectively free.
\end{itemize}
A zero-capital $\adv$ cannot produce a valid \PoN{} block, cannot claim a tier it has not
collateralised (\textbf{A4}), and cannot reach any collateral-weighted threshold.

\paragraph{What capital buys.} Write $\vtotal = 88{,}514.5$ for the measured total
collateral weight, in the units of $1{,}000\flux$ of \Cref{tab:notation} (that is,
$88{,}514{,}500\flux$ of collateral), and $\ncount{\mathrm{tot}} = 6{,}489$ for the measured node count
\srcmeasured{api.runonflux.io/\allowbreak{}daemon/\allowbreak{}listzelnodes}{2026-09-03}.

\begin{property}[Share of block production]
\label{prop:share}
\shipped Under \textbf{A1} and \textbf{A3}, let $\adv$ control $m = f\,\ncount{\mathrm{tot}}$
confirmed nodes. Then $\adv$'s expected share of produced blocks is at least $f$, and strictly
exceeds $f$ whenever the per-slot eligibility probability $p$ is bounded away from $0$.
\end{property}

\begin{proof}
Eligibility in slot $\tslot$ is $H(\text{collateral}\Vert\text{prevhash}\Vert\tslot) \le T$, one hash
per confirmed node, \emph{independent of tier}
\srccode{fluxd/\allowbreak{}src/\allowbreak{}pon/\allowbreak{}pon.cpp}{70--93,400--433}. Distinct collateral outpoints give distinct
preimages, so per slot eligibility is Bernoulli($p$) independently across nodes. A slot yields a
block iff at least one node is eligible, with probability $g(\ncount{\mathrm{tot}})$ where
$g(x) = 1-(1-p)^x$. $\adv$ takes the slot whenever at least one of its $m$ nodes is eligible: the
rank-based delay of $\qty{4000}{\milli\second}$ per rank is documented in the source as an
orphan-reduction heuristic with no consensus enforcement
\srccode{fluxd/\allowbreak{}src/\allowbreak{}pon/\allowbreak{}pon-minter.cpp}{78--104,217--221,250}, so an $\adv$ that ignores it publishes
before any honest node of rank $\ge 1$; the \qty{5}{\second} minimum-since-last-block guard in the
same minter loop is likewise unenforced by consensus \srccode{fluxd/\allowbreak{}src/\allowbreak{}pon/\allowbreak{}pon-minter.cpp}{200--207},
and the only consensus bound is that a block's time exceed its predecessor's
\srccode{fluxd/\allowbreak{}src/\allowbreak{}main.cpp}{5564--5572}. Hence
$\adv$'s share is $g(m)/g(\ncount{\mathrm{tot}})$. The map $g$ is concave on $[0,\infty)$ with
$g(0)=0$, so $g(f\ncount{\mathrm{tot}}) \ge f\,g(\ncount{\mathrm{tot}}) + (1-f)g(0) =
f\,g(\ncount{\mathrm{tot}})$, giving share $\ge f$; the inequality is strict for $p>0$ by strict
concavity. The per-slot model understates $\adv$: consensus does not pin a block's slot to
wall-clock time, so $\adv$ may evaluate eligibility over every slot in the \qty{300}{\second}
future-time window rather than the current slot alone (\S\ref{sec:consensus-sortition}), which can
only raise its share, so the bound stands.
\end{proof}

\begin{property}[Cost of an eligibility share]
\label{prop:eligibility-cost}
\shipped To reach a fraction $f$ of eligibility draws by newly collateralising nodes,
$\adv$ must add $m = \frac{f}{1-f}\ncount{\mathrm{tot}}$ nodes, at minimum cost
$\advbudget = m\cdot\coll{\tC}$. At $f = 1/2$ this is $m = 6{,}489$ nodes and
$\advbudget = 6{,}489{,}000\flux$. Per eligibility draw, $\tC$ is exactly
$\coll{\tS}/\coll{\tC} = 40$ times cheaper than $\tS$.
\end{property}

\begin{proof}
$m/(\ncount{\mathrm{tot}}+m) = f$ rearranges to the stated $m$. Cost is minimised by the cheapest
tier because, by the proof of Property~\ref{prop:share}, sortition assigns one hash per confirmed
node regardless of tier. The registration is visible on-chain in advance: a start transaction
requires $100$ blocks of collateral maturity \srccode{fluxd/\allowbreak{}src/\allowbreak{}fluxnode/\allowbreak{}fluxnode.h}{20} and the
confirm transaction is a separate on-chain event.
\end{proof}

\begin{property}[Cost of the moderation threshold]
\label{prop:moderation-cost}
Let $\vthresh$ be a removal threshold expressed as a fraction of collateral weight $\vtotal$
(\S\ref{sec:governance}, \planned). Then, in $\vtotal$'s units of $1{,}000\flux$,
(i) acquiring existing collateral costs $\advbudget = \vthresh\,\vtotal$, because the transfer leaves
$\vtotal$ unchanged; (ii) newly collateralising costs
$\advbudget = \frac{\vthresh}{1-\vthresh}\,\vtotal$. At $\vthresh = 0.25$ these are
$22{,}128{,}625\flux$ and $29{,}504{,}833\flux$ respectively.
\end{property}

\begin{proof}
(i) is immediate. For (ii), solve $w/(\vtotal+w) = \vthresh$. The measured distribution corroborates
(i): the three largest node keys hold $22{,}200{,}000\flux$, which is $25.08\%$ of $\vtotal$
\srcmeasured{api.runonflux.io/\allowbreak{}daemon/\allowbreak{}listzelnodes}{2026-09-03} --- four addresses under the
owner-identity model (\S\ref{sec:governance}; \textbf{A10}) --- so an $\adv$ purchasing that same
weight reaches the threshold, and does so by transacting with three or four counterparties.
\end{proof}

\begin{remark}[The two mechanisms are not independently priced]
\label{rem:cross-price}
Composing Properties~\ref{prop:eligibility-cost} and~\ref{prop:moderation-cost}: an $\adv$ that
reaches $\vthresh = 0.25$ by newly collateralising $29{,}504$ $\tC$ nodes simultaneously holds
$29{,}504/(6{,}489+29{,}504) = 82.0\%$ of eligibility draws. Conversely the $6{,}489{,}000\flux$ that
buys half of all eligibility draws buys only $6{,}489/(88{,}514.5+6{,}489) = 6.8\%$ of moderation
weight. The cheapest route to block production is the most expensive route to votes and vice versa,
and the cheap route to the moderation threshold is also a route to block-production dominance. Both
routes are visible on-chain in advance, which is the property \S\ref{sec:governance} builds on.
\end{remark}

\begin{property}[Reorganisation cost]
\label{prop:reorg}
\shipped Under \textbf{A1} and \textbf{A3}, the cost of an $\ell$-block reorganisation is a function
of collateral and of elapsed slots, not of computation, and the capital committed is recoverable.
\end{property}

\begin{proof}[Proof sketch]
Every \PoN{} block contributes a fixed $2^{40}$ to chainwork regardless of its target
\srccode{fluxd/\allowbreak{}src/\allowbreak{}pow.cpp}{284--290}, so a private chain outranks the public chain exactly when its
block count is larger. By the proof of Property~\ref{prop:share}, an $\adv$ holding fraction $f$ of
confirmed nodes extends its private chain at expected rate $g(f\ncount{\mathrm{tot}})$ blocks per
slot, with $g(x)=1-(1-p)^x$; producing $\ell$ private blocks therefore costs $\adv$ elapsed slots,
not hash trials, and the required capital is $\advbudget$ from
Property~\ref{prop:eligibility-cost}. That capital is not consumed by the attempt: the only three
transitions that remove a node from the paid set are collateral spent, failure to re-confirm, and a
migration-window amount mismatch (Definition~\ref{def:operator}), none of which destroys value. The
binding constraint on $\ell$ is therefore the $40$-block reorganisation cap
\srccode{fluxd/\allowbreak{}src/\allowbreak{}main.h}{74}, not the cost of the attempt.
\end{proof}

\begin{remark}[Fee-pool income]
\label{rem:pool-share}
\planned Under \PNR{} (\S\ref{sec:pnr}) a node would additionally receive a share of the fee pool
$\pool$ once per rotation $\rot{t}$ of its tier --- at the measured population, once every $1{,}627$
blocks for $\tS$, which is $\qty{13.56}{\hour}$ at $\spacingpon$
\srcmeasured{api.runonflux.io/\allowbreak{}daemon/\allowbreak{}listzelnodes}{2026-09-03}. The magnitude of that share is not
fixed here.
\settlednote{$\adv$'s per-node fee-pool income under \PNR{} is
$\tierratio{t}\,D_\hgt/\ncount{t}$ per block on average, with $D_\hgt = \lfloor \pool_{\hgt-1}/\drip \rfloor$
(Proposition~\ref{prop:pnr:across}) --- the per-block drip, times the tier's split, divided by the
confirmed nodes in that tier; the operator share $\nodeshare$ enters through the pool's inflow rather
than at payout. Every term is now fixed:
the drip is $\pool_{\hgt-1}/K$ with $K = \num{86400}$, the operator share follows the ramp of
\S\ref{sec:pnr} to its cap, and the tier split is the settled partition
(\Cref{tab:pnr-recommended}). This section still assumes no \emph{value} for it, because the pool
balance is demand-determined and the adversary model must hold for any balance --- but the
\emph{function} is no longer open.}
\end{remark}

\subsection{Assumptions}
\label{sec:model:assumptions}

Each assumption is cited elsewhere by its label. Each states what it enables and what it forecloses.

\begin{assumption}[\textbf{A1} --- chain synchrony]
\label{asm:a1-synchrony}
\shipped The chain network is partially synchronous in the sense of Definition~\ref{def:chainnet},
and \PoN{} additionally assumes clock skew bounded by the \qty{300}{\second} header tolerance
\srccode{fluxd/\allowbreak{}src/\allowbreak{}main.cpp}{5330--5378}.
\emph{Enables}: liveness arguments for block production, and settlement arguments at depth.
\emph{Forecloses}: any argument assuming synchrony, and any use of median-time-past ordering for
\PoN{} blocks, which use strict predecessor-time ordering instead
\srccode{fluxd/\allowbreak{}src/\allowbreak{}main.cpp}{5564--5572}.
\end{assumption}

\begin{assumption}[\textbf{A2} --- overlay delivery]
\label{asm:a2-overlay}
\shipped The orchestration overlay delivers authenticated messages best-effort, with convergence
conditional on the peer thresholds and TTLs of Definition~\ref{def:overlay}
\srccode{flux/\allowbreak{}ZelBack/\allowbreak{}config/\allowbreak{}default.js}{262--271,311--314}.
\emph{Enables}: eventual-convergence statements about placement, conditional on peer counts.
\emph{Forecloses}: any safety property over application state, any bound on time to convergence, and
any chain-derived statement about placement (Property~\ref{prop:two-models}).
\end{assumption}

\begin{assumption}[\textbf{A3} --- fork choice is a block-count race]
\label{asm:a3-forkchoice}
\shipped Every \PoN{} block contributes a fixed chainwork of $2^{40}$, independent of its difficulty
target \srccode{fluxd/\allowbreak{}src/\allowbreak{}pow.cpp}{284--290}. Fork choice therefore compares block counts, resolved
by first-seen, and the rank-based minting delay is advisory
\srccode{fluxd/\allowbreak{}src/\allowbreak{}pon/\allowbreak{}pon-minter.cpp}{78--104,217--221}.
\emph{Enables}: reorganisation cost derived from collateral and eligibility
(Property~\ref{prop:reorg}); a security argument that does not depend on hashing.
\emph{Forecloses}: every cumulative-work argument, including light-client fork choice
(\textbf{A9}) and any claim that reorganisation cost rises with hash rate.
\notestablished{\texttt{n\allowbreak{}Minimum\allowbreak{}Chain\allowbreak{}Work} still carries its pre-\PoN{} value
\srccode{fluxd/\allowbreak{}src/\allowbreak{}chainparams.cpp}{168--170}; whether it remains an effective initial-block-download
guard under fixed-work \PoN{} was not resolved.}
\end{assumption}

\begin{assumption}[\textbf{A4} --- benchmark results are collateral-bound, not hardware-attested]
\label{asm:a4-benchmark}
\shipped A node's tier claim is carried in a confirmation transaction signed by \fluxbench{} using
keys shared across the fleet, validated by consensus against a maintainer allowlist
\srccode{fluxd/\allowbreak{}src/\allowbreak{}chainparams.cpp}{233--240}, and obtained by \fluxd{} from the local
\texttt{fluxbench-cli} through \texttt{popen()}, with \texttt{status:\allowbreak{}"complete"} treated as
sufficient \srccode{fluxd/\allowbreak{}src/\allowbreak{}fluxnode/\allowbreak{}benchmarks.cpp}{227--291}. The registration must match an
on-chain collateral amount for the claimed tier \srccode{fluxd/\allowbreak{}src/\allowbreak{}coins.cpp}{630--686}.
\emph{Enables}: the statement that tier is collateral-bound --- an operator cannot claim a tier it
has not collateralised. \emph{Forecloses}: any property conditioned on the host actually meeting the
tier's hardware thresholds, and any use of a benchmark result as evidence about a specific machine.
This is a benchmark-gating bypass, not tier forgery, and \S\ref{sec:consensus} and
\S\ref{sec:attestation} must both be written against the weaker statement.
\end{assumption}

\begin{assumption}[\textbf{A5} --- attestations are not machine-bound]
\label{asm:a5-attestation}
\shipped A \fluxsas{} attestation demonstrates possession of a value present in every copy of the
software; none of the derivation inputs depends on the TPM, on Secure Boot state, or on any
per-machine value \srcprivate{\SAS{} API server}. Every
\texttt{/v2/*} route requires an internal client certificate and is bound to loopback, so no
tenant-facing verification path exists \srccode{fluxsas/\allowbreak{}src/\allowbreak{}api\_server.h}{1042--1053,220--252}.
\emph{Enables}: measured boot, dm-verity and LUKS as genuine host-local defences against a remote
attacker. \emph{Forecloses}: any protocol that consumes an attestation as a machine binding, any
tenant-verifiable execution claim, and any argument that the eligibility gate is enforced by
hardware. The bonded attestation network of \S\ref{sec:attestation} is the answer to exactly this.
\end{assumption}

\begin{assumption}[\textbf{A6} --- emergency blocks bypass \PoN{} unconditionally]
\label{asm:a6-emergency}
\shipped A block whose collateral field equals the emergency sentinel and which carries $2$ valid
signatures from a hardcoded set of $4$ public keys is accepted, bypassing \texttt{Check\allowbreak{}Proof\allowbreak{}Of\allowbreak{}Node}
and per-node signature verification \srccode{fluxd/\allowbreak{}src/\allowbreak{}emergencyblock.cpp}{56--150},
\srccode{fluxd/\allowbreak{}src/\allowbreak{}pon/\allowbreak{}pon.cpp}{205--252}. The only gate is that \PoN{} is active: there is no time
delay, no cooldown and no stall detection, despite an in-source comment stating that such
restrictions were contemplated \srccode{fluxd/\allowbreak{}src/\allowbreak{}emergencyblock.cpp}{183--191}.
\emph{Enables}: recovery from a chain halt without a coordinated software release.
\emph{Forecloses}: any safety property quantified over all valid blocks under \PoN{} eligibility; any
light-client rule derived from eligibility (\textbf{A9}); and, as a forward constraint,
\S\ref{sec:pnr} must exclude emergency blocks from the fee-pool drip, or the emergency key set,
facing no time gating, inherits the ability to accelerate the release of accumulated fees. It could
not redirect them: an emergency block still fills the deterministic node payout
\srccode{fluxd/\allowbreak{}src/\allowbreak{}miner.cpp}{487--489}, and connection rejects any block that omits it
\srccode{fluxd/\allowbreak{}src/\allowbreak{}main.cpp}{3878--3884}.

\emph{Planned change.} The team intends to move this set to $3$-of-$4$ with newly generated keys,
aligning it with the bridge quorum of \S\ref{sec:bridging} \planned
\srcintent{design intake}. No activation is scheduled, and the assumption is
unchanged in kind: at $3$-of-$4$ the path still bypasses \PoN{} eligibility and per-node signature
verification with no time or stall-detection gating, so every property that references A6 continues
to hold.
\end{assumption}

\begin{assumption}[\textbf{A7} --- discovery, naming and runtime policy are operated]
\label{asm:a7-operated}
\shipped Application discovery has a single source \srccode{flux-domain-manager/\allowbreak{}src/\allowbreak{}services/\allowbreak{}flux/\allowbreak{}dataFetcher.js}{47--91};
DNS, HAProxy election and certificate issuance run on a $16$-host fleet inside one private
\texttt{/24} \srccode{flux-fdm-infrastructure/\allowbreak{}ansible/\allowbreak{}inventory/\allowbreak{}hosts.yaml}{1--206}; and the runtime
network policy is a set of unsigned documents fetched over plain HTTPS and enforced by every node
\srccode{flux/\allowbreak{}ZelBack/\allowbreak{}src/\allowbreak{}services/\allowbreak{}appSecurity/\allowbreak{}imageManager.js}{179--195}, whose authenticity
``rests on GitHub and on who can merge'' and which takes fleet-wide effect on the next scheduled
fetch of \qty{6}{\hour}--\qty{12}{\hour}, with no staged rollout and no second-reviewer requirement
\srcdoc{fluxos-network-policy/README.md:30--44,279--296,320--326}.
\emph{Enables}: incident response measured in hours; the honest baseline against which
\S\ref{sec:governance} argues that a collateral-weighted quorum is an improvement.
\emph{Forecloses}: any end-to-end decentralization claim for reachability, naming or content policy,
and any censorship-resistance claim about application availability as deployed.
\end{assumption}

\begin{assumption}[\textbf{A8} --- the shielded pool is closed to new entrants]
\label{asm:a8-shielded}
\shipped Since height $835{,}554$, consensus rejects any transaction carrying transparent inputs
together with a joinsplit or a shielded output, with reject reason
\texttt{bad-\allowbreak{}txns-\allowbreak{}fluxrebrand-\allowbreak{}vprivacy-\allowbreak{}not-\allowbreak{}empty} \srccode{fluxd/\allowbreak{}src/\allowbreak{}main.cpp}{1398--1408}. Shielded
value may move from one \zaddr to another, or exit from a \zaddr to a \taddr; it cannot enter.
\emph{Enables}: an exact statement of the deployed privacy surface, and the measurement that
$96{,}479.94\flux$ sits in the two pools
\srcmeasured{api.runonflux.io/\allowbreak{}daemon/\allowbreak{}getblockchaininfo}{2026-09-03}.
\emph{Forecloses}: any construction in \S\ref{sec:private-renewal} that requires funding a shielded
note from transparent value, and any claim that the anonymity set grows. Re-opening the pool would be
a consensus change; the team has instead decided to retire both pools
(\S\ref{sec:shielded-retirement}, \srcintent{design intake}), so this assumption
describes a mechanism on its way out rather than a prerequisite for future work.
\end{assumption}

\begin{assumption}[\textbf{A9} --- no light client is constructible]
\label{asm:a9-lightclient}
\shipped With the current header format no \Flux{} light client can be constructed, for three
compounding reasons: the \PoN{} header carries no commitment to the node set, so eligibility cannot
be checked without the full registry \srccode{fluxd/\allowbreak{}src/\allowbreak{}primitives/\allowbreak{}block.h}{44--45,64--96}; chainwork
is fixed per block (\textbf{A3}); and the emergency path can produce a valid block outside the
eligibility rules (\textbf{A6}).
\emph{Enables}: an honest enumeration of the bridge design space in \S\ref{sec:bridge}.
\emph{Forecloses}: light-client bridging and any trust-minimized verification of \Flux{} state on
another chain, until a header commitment to the node set exists (\S\ref{sec:chain-ops}).
\end{assumption}

\begin{assumption}[\textbf{A10} --- identity is collateral and nothing else]
\label{asm:a10-identity}
\shipped There is no Sybil-resistant identity in the system other than a collateralised node. The two public
proxies for operator identity, \texttt{pubkey} and \texttt{payment\_address}, are not proof of common
ownership in either direction, so every concentration figure in this paper is a \emph{lower bound} on
concentration \srcmeasured{api.runonflux.io/\allowbreak{}daemon/\allowbreak{}listzelnodes}{2026-09-03}.
\emph{Enables}: Sybil resistance by construction for any collateral-weighted mechanism.
\emph{Forecloses}: honest-majority-of-identities assumptions, one-node-one-vote designs, and any
claim derived from counting distinct keys as if they were distinct parties.
\end{assumption}

\begin{assumption}[\textbf{A11} --- denomination and exogeneity of price]
\label{asm:a11-denomination}
All economic quantities in this paper are denominated in \flux{} or in resource units. The market
price of \flux{} is exogenous and is never modelled, forecast or used in any derivation.
\emph{Enables}: statements about issuance, supply, operator income and demand that a reader can check
against the chain. \emph{Forecloses}: any fiat-denominated claim, and any crossover or parity result
that would depend on price rather than on demand (\S\ref{sec:pnr}, \S\ref{sec:evaluation}).
\end{assumption}

\subsection{Out of scope}
\label{sec:model:scope}

Three things are outside the model and are never assumed, argued or defended.
\begin{enumerate}
  \item \emph{The security of the container runtime.} Nothing in the attestation path measures or
        attests the Docker daemon, \FluxOS{} itself, or the hypervisor
        \srccode{fluxsas/\allowbreak{}src/\allowbreak{}watchdog.h}{547--559}. Isolation between co-tenant containers on one
        host is assumed to be whatever the runtime provides, and no property in this paper depends on
        it being more than that.
  \item \emph{The correctness of tenant workloads.} A tenant's image may do anything the runtime
        permits. Content policy is treated only where it is a network mechanism
        (\S\ref{sec:governance}); workload correctness is the tenant's.
  \item \emph{The market price of \flux{}} (\textbf{A11}).
\end{enumerate}

\subsection{Notation}
\label{sec:model:notation}

Table~\ref{tab:notation} lists the symbols used in this paper, grouped by domain, with the domain or
unit of each. The full index, including every constant and its provenance, is Appendix~D; the
constant tables are Appendix~A. Macros are defined once in \texttt{notation.tex} and are not
redefined anywhere in the text.

\begin{table}[H]
\centering
\small
\caption{Notation by domain. Full index: Appendix~D.}
\label{tab:notation}
\begin{tabular}{@{}p{0.14\linewidth}p{0.50\linewidth}p{0.30\linewidth}@{}}
\toprule
\textbf{Symbol} & \textbf{Meaning} & \textbf{Domain / unit} \\
\midrule
\multicolumn{3}{@{}l}{\textit{Chain and consensus}}\\
$\hgt$ & block height & $\mathbb{N}$ \\
$\hpon$, $\hpa$ & \PoN{} activation height; PA-depletion height & $\mathbb{N}$ \\
$\tslot$ & \PoN{} slot index & $\mathbb{N}$ \\
$\spacingpon$, $\spacingpow$ & \PoN{} / legacy target spacing & \unit{\second} \\
$\blocksyear$ & blocks per year & blocks \\
$\redint$, $\redmax$, $\redfac$ & subsidy reduction interval, count, factor & blocks; $\mathbb{N}$; $\mathbb{Q}\cap(0,1)$ \\
\midrule
\multicolumn{3}{@{}l}{\textit{Tiers}}\\
$\tiers$ & $\{\tC,\tN,\tS\}$ & finite set \\
$\coll{t}$ & collateral required by tier & $\tiers \to \mathbb{Q}_{>0}$, \flux \\
$\ncount{t}$, $\rot{t}$ & node count; rotation length & $\mathbb{N}$; blocks \\
$\tbase{t}$ & tier base against a total of $14$ & $\tiers \to \mathbb{Q}_{>0}$, \flux \\
\midrule
\multicolumn{3}{@{}l}{\textit{Emission and supply}}\\
$\subsidy$, $\subsidypon$, $\subsidypow$ & block subsidy & $\mathbb{N}\to\mathbb{Q}_{\ge0}$, \flux \\
$\devfund$ & consensus dev-fund remainder & \flux{} per block \\
$\supply$, $\premine$ & cumulative issued supply; premine & \flux \\
\maxmoney & per-transaction sanity bound, \emph{not} a cap & \flux \\
$\sat$ & $10^{-8}\,$\flux & unit \\
\midrule
\multicolumn{3}{@{}l}{\textit{Applications and pricing}}\\
$\app$, $\ninst$, $\dur$ & an application; instance count; duration & --- ; $\mathbb{N}$; blocks \\
$\resvec$, $\pricevec$ & resource vector; price vector at a height & $\mathbb{Q}_{\ge0}^{3}$ \\
$\price$ & price of a deployment & \flux \\
$\blockslasting$ & default lifetime & blocks \\
\midrule
\multicolumn{3}{@{}l}{\textit{\PNR{} (\planned, \S\ref{sec:pnr})}}\\
$\pool$, $\drip$ & fee-pool balance; drip constant & \flux; blocks \\
$\nodeshare$, $\nodesharecap$, $\rampstep$ & operator share, ceiling, per-step increase & $[0,1]$ \\
$\payees$, $\tierratio{t}$ & nodes paid in a block; tier split of the drip & finite set; $[0,1]$ \\
\midrule
\multicolumn{3}{@{}l}{\textit{Moderation (\planned, \S\ref{sec:governance})}}\\
$\vweight{i}$, $\vtotal$ & voting weight of identity $i$; total weight & units of $1{,}000\flux$ \\
$\vthresh$, $\bond$, $\reporters$ & removal threshold; report bond; reporters & $[0,1]$; \flux; finite set \\
\midrule
\multicolumn{3}{@{}l}{\textit{Privacy and value}}\\
\taddr, \zaddr & transparent / shielded address & --- \\
$\note$, $\nf$, $\cm$, $\tree$ & note, nullifier, commitment, commitment tree & --- \\
$\zkproof$, $\viewkey$ & zero-knowledge proof; incoming viewing key & --- \\
\midrule
\multicolumn{3}{@{}l}{\textit{Agents and delegation (\planned, \S\ref{sec:agent-identity})}}\\
$\principal$, $\agent$ & human principal; autonomous agent & --- \\
$\budget$, $\scopeset$ & delegated spending budget; scope of authority & \flux; finite set \\
\midrule
\multicolumn{3}{@{}l}{\textit{Adversary}}\\
$\adv$, $\advbudget$ & the adversary; its capital budget & --- ; \flux \\
\bottomrule
\end{tabular}
\end{table}

\subsection{How to read the tags}
\label{sec:model:tags}

Every substantive claim in this paper carries one \emph{maturity} tag and one \emph{provenance} tag.
This is a promise to the reader, and it is the mechanism by which a reader can tell, at any point,
which half of the paper they are in.

\paragraph{Maturity.} \shipped means the behaviour is in code that is deployed and reachable on
mainnet. \inprogress means the code exists and is under active development but is not activated on
any network. \planned means a design exists and is specified in this paper but no implementation is
deployed. \vision means a direction with requirements and interfaces but no construction.

\paragraph{Provenance.} \srccode{repo/\allowbreak{}file}{line} cites deployed source. \srcbranch{repo}{branch}{file:line}
cites a branch, which a reader may not be able to fetch. \srcdoc{name}, \srcroadmap{item} and
\srcintent{file}
cite documents, and are hearsay about behaviour. \srcmeasured{endpoint}{date} cites a public endpoint
with its retrieval date. \srcinferred{} marks a conclusion drawn by the author from cited material
rather than stated in it; it never supports \shipped. Citations into \texttt{flux/} are against the
corpus checkout, commit \texttt{933614b1c} on \texttt{feat/\allowbreak{}appspec-\allowbreak{}v9}; in
\texttt{ZelBack/\allowbreak{}config/\allowbreak{}default.js} the line numbers above 127 differ from
\texttt{master} by up to fifteen lines.

\paragraph{The rules.}
\begin{enumerate}
  \item \shipped requires a source citation or a measurement citation. A document alone is never
        sufficient to support \shipped; a roadmap or a README describes intent, not behaviour.
  \item \planned and \vision material never appears in the present tense. Where this paper writes
        ``the pool would pay'' rather than ``the pool pays'', that is not hedging --- it is the tag
        being honoured in the grammar.
  \item A property the author cannot argue appears as a \texttt{Conjecture}, an
        \texttt{Open Problem}, or a \emph{Not established by this survey} callout. It never appears
        as an unlabelled assertion.
  \item An undecided parameter carries its domain, the trade-off, and --- once decided --- a
        \emph{Settled} callout recording the value. A guess is never inserted in place of a decision.
\end{enumerate}

%% Drain this section's float queue before the next begins.
\FloatBarrier
